% 由 tools/bench/analysis/render_tables.py 生成，请勿手改。
% 直接由 contracts/golden-vectors/cfsi_stress_pca_q90_reference.json 计算。
\begin{table}[t]
\centering
\caption{标签循环性诊断（$y_{20}$ 为 $t+1$ 至 $t+20$ 的严格未来标签；\texttt{strict} 再推后一天至 $t+2$）}
\label{tab:circularity}
\footnotesize
\begin{tabular}{lr}
\toprule
统计量 & 值 \\
\midrule
样本日数（$y_{20}$ 可用） & 2646 \\
基准发生率 $P(y_{20}=1)$ & 0.341 \\
$P(y_{20}=1\mid hs_t=1)$ & 0.921 \\
$P(y_{20}=1\mid hs_t=0)$ & 0.265 \\
严格标签样本日数 & 2646 \\
$y_{20}$ 与严格标签一致的日数 & 2646 \\
基准发生率 $P(y^{\mathrm{strict}}=1)$ & 0.341 \\
$P(y^{\mathrm{strict}}=1\mid hs_t=1)$ & 0.921 \\
$P(y^{\mathrm{strict}}=1\mid hs_t=0)$ & 0.265 \\
\bottomrule
\end{tabular}
\par\smallskip
\raggedright\scriptsize
注：$hs_t$ 为当日 high\_stress（定义即 CFSI $>\tau$）。
随附黄金向量中 $y_{20}$ 与 $t+1$ 至 $t+20$ 的严格未来标签逐日一致（见上行），
即本产物口径不含当日（与 \texttt{docs/13} §5.5 旧表述不一致，差异见附录 A.4）。
无论口径如何，二者都与 CFSI 同源，
$P(y\mid hs_t=1)$ 只能作为循环性诊断，不能当作预测能力。
\end{table}
